\documentclass[11pt,a4paper]{extarticle}
\pagenumbering{gobble}
\usepackage[left=0.7in,right=0.7in,top=0.5in,bottom=0.5in]{geometry}
% XeLaTeX ile derlenir: xelatex -jobname=enesorhan_resume_tr_final enesorhan_resume_tr_duzenlenmis.tex
% Gerçek Türkçe glifleri (ğ, ş, ı, İ) sayesinde metin ATS tarafından doğru okunur.
\usepackage{fontspec}
\setmainfont{Helvetica Neue}
\usepackage[turkish,shorthands=off]{babel}
\usepackage{xcolor}
\usepackage{titlesec}
\usepackage{enumitem}
\usepackage{array}
\usepackage[unicode]{hyperref}


\definecolor{accent}{HTML}{1A365D}
\definecolor{muted}{HTML}{4A5568}
\definecolor{rulecol}{HTML}{B8C4D4}

\hypersetup{colorlinks=true,urlcolor=accent,linkcolor=accent,
            pdftitle={Enes Orhan - Backend Developer},pdfauthor={Enes Orhan}}

%% ---- Gövde yazısı: 10,5 pt, sola hizalı, kelime bölme yok ---------------
\AtBeginDocument{%
  \fontsize{10.5}{13.6}\selectfont
  \raggedright
  \hyphenpenalty=10000 \exhyphenpenalty=10000
}
\setlength{\parindent}{0pt}
\setlength{\parskip}{0pt}

%% ---- Bölüm başlıkları ----------------------------------------------------
\titleformat{\section}
  {\large\bfseries\color{accent}}{}{0em}{}
  [\vspace{-2pt}{\color{rulecol}\titlerule[0.8pt]}]
\titlespacing*{\section}{0pt}{10pt}{5pt}

%% ---- Yapı taşları --------------------------------------------------------
\newcommand{\sep}{\ \textperiodcentered\ }
\newcommand{\ca}{\textasciitilde}

% Bölünmeyen blok: bir deneyim/proje sayfalar arasında ikiye ayrılmaz
\newenvironment{block}[1][]
  {\par\noindent\begin{minipage}[t]{\linewidth}\raggedright
   \hyphenpenalty=10000 \exhyphenpenalty=10000
   \if\relax\detokenize{#1}\relax\else\vspace{4pt}\section{#1}\fi}
  {\end{minipage}\par\vspace{7pt}}

% Başlık satırı: ad (kalın) + sağda tarih
\newcommand{\head}[2]{{\bfseries #1}\hfill{\color{muted}#2}\par}
% Alt satır: rol, konum, bağlantılar
\newcommand{\sub}[1]{{\color{muted}\itshape #1}\par}
% Bağlam cümlesi
\newcommand{\ctx}[1]{\vspace{2pt}{\color{muted}#1}\par}
% Teknoloji satırı
\newcommand{\tech}[1]{\vspace{2pt}{\small\color{muted}#1}\par}

\setlist[itemize]{leftmargin=1.1em,label=\textbullet,itemsep=2pt,topsep=3pt,parsep=0pt,partopsep=0pt}

\begin{document}

%% ======================= BAŞLIK =======================
\begin{center}
{\fontsize{22}{26}\selectfont\bfseries\color{accent}Enes Orhan}\\[5pt]
{\large\color{muted}Backend Developer\;{\color{rulecol}\textbar}\;C\#, ASP.NET Core, Django}\\[6pt]
Adana, Türkiye\sep
\href{tel:+905077561488}{+90 507 756 14 88}\sep
\href{mailto:enesorhann404@gmail.com}{enesorhann404@gmail.com}\\[2pt]
\href{https://linkedin.com/in/enesorhann}{linkedin.com/in/enesorhann}\sep
\href{https://github.com/enesorhann}{github.com/enesorhann}\sep
\href{https://enesorhan.com}{enesorhan.com}
\end{center}

%% ======================= ÖZET =======================
\section{Özet}
2026'da Konya Teknik Üniversitesi Bilgisayar Mühendisliği bölümünden mezun oldum ve backend geliştirmeye odaklanıyorum.
C\# ve ASP.NET Core ile App Store ve Google Play'de yayında olan, gerçek kullanıcıları bulunan iki uygulamanın backend'ini geliştirdim.
Profesyonel deneyimimi Django ile edindim: bir fiyat karşılaştırma platformunda üç kişilik ekipte, ardından bir B2B ihale
platformunda backend'in ana geliştiricisi olarak çalıştım. Projelerimi Git ile yönetiyor, Docker ve Kubernetes ile yayına alıyorum.

%% ======================= .NET PROJELERİ =======================
\begin{block}[Yayındaki ASP.NET Core Projeleri]
\head{Zerbook}{Mart 2026 -- Halen}
\sub{Kişisel ürün\sep Altın birikimi takibi\sep 36 kayıtlı kullanıcı\sep
\href{https://apps.apple.com/tr/app/zerbook-alt\%C4\%B1n-takip-portf\%C3\%B6y/id6800002144?l=tr}{App Store} /
\href{https://play.google.com/store/apps/details?id=com.zerbook.zerbook_mobile}{Google Play}}
\begin{itemize}
\item Kullanıcıların altın birikimlerini kaydedip güncel fiyatla değerini gördüğü REST API'yi katmanlı bir mimariyle geliştirdim.
\item Ailelerin birikimlerini birlikte takip etmesi için davet kodlu gruplar, üye rolleri ve üyeler arası transfer ekledim.
\item Miktar ve fiyatları tamsayı birimlerde (mikrogram, adet, kuruş) tutarak kayan nokta kaynaklı hassasiyet hatalarını önledim.
\item Fiyatları dakikada bir güncelleyen arka plan servisini kurdum. E-posta, telefon, Google ve Apple girişlerini
      Firebase üzerinden API'nin JWT oturumuna bağladım.
\end{itemize}
\tech{C\#\sep ASP.NET Core\sep EF Core\sep PostgreSQL\sep JWT\sep Firebase\sep xUnit\sep Docker\sep Kubernetes}
\end{block}

\begin{block}
\head{Entegram}{Nisan 2026 -- Halen}
\sub{Kişisel ürün\sep İngilizce öğrenme uygulaması\sep 57 kayıtlı kullanıcı\sep
\href{https://apps.apple.com/tr/app/entegram-kayd\%C4\%B1r-\%C3\%B6\%C4\%9Fren/id6766254758?l=tr}{App Store} /
\href{https://play.google.com/store/apps/details?id=com.enesorhan.gibigram.gibigram_mobile}{Google Play}}
\begin{itemize}
\item Kullanıcıların eş zamanlı yarıştığı quiz düellolarını ve anlık mesajlaşmayı SignalR ile geliştirdim.
\item Kaydedilen kelimeleri SM-2 algoritmasıyla tekrar zamanına göre planlayan sistemi test edilebilir biçimde yazdım.
\item iOS ve Android aboneliklerini RevenueCat ile backend'de eşitledim; premium erişimi sunucu tarafında kontrol ediyorum.
\item Bildirim kampanyaları gibi zamanlanmış işleri, sunucu yeniden başlasa da kaybolmayacak şekilde Hangfire ile çalıştırdım.
\end{itemize}
\tech{C\#\sep ASP.NET Core\sep SignalR\sep Hangfire\sep PostgreSQL\sep RevenueCat\sep Docker\sep Kubernetes}
\end{block}

%% ======================= DENEYİM =======================
\begin{block}[Deneyim]
\head{BidsForShips\sep\href{https://bidsforships.com}{bidsforships.com}}{Ekim 2025 -- Şubat 2026}
\sub{Full Stack Developer (Freelance)\sep Uzaktan}
\ctx{B2B denizcilik ihale platformu. Kurucunun hazırladığı Django altyapısı üzerinde backend'in büyük bölümünü geliştirdim;
frontend'de bir geliştiriciyle birlikte çalıştım.}
\begin{itemize}
\item Teklif isteme, teklif verme, kısa liste, kabul ve imzalı sözleşme aşamalarından oluşan ihale sürecini
      rol bazlı yetki ve NDA kontrolleriyle geliştirdim.
\item Teklif ve ekipman revizyonlarının geçmişini saklayan yapıyı kurdum; taraflar belgelerdeki değişiklikleri takip edebiliyor.
\item Ana ve alt şirket yapısını, şirket yöneticisi davetini ve tedarikçi malzeme talebi modülünü geliştirdim;
      3.000+ şirket kaydını Excel'den aktaran komutu yazdım.
\item E-posta bildirimlerini Celery ile arka plana taşıdım; Next.js ve TypeScript ile panel ve teklif ekranlarını geliştirdim.
\end{itemize}
\ctx{Kaynak kodlar Fiyator’un şirket GitHub organizasyonundaki özel reposunda bulunuyor.}
\tech{Python\sep Django REST Framework\sep PostgreSQL\sep Celery\sep Elasticsearch\sep Next.js\sep TypeScript\sep Kubernetes}
\end{block}

\begin{block}
\head{Fiyator A.Ş.\ \textmd{(Ankara)}}{Haziran 2025 -- Ekim 2025}
\sub{Software Engineer Intern, ardından Freelance\sep Uzaktan}
\ctx{14M+ ürün kaydı olan fiyat karşılaştırma platformu. Kurucu, bir mobil geliştirici ve benden oluşan ekipte çalıştım.}
\vspace{3pt}
{\itshape Freelance}\hfill{\color{muted}Ağustos 2025 -- Ekim 2025}\par
\begin{itemize}
\item ikas, Ticimax ve XML kaynaklarından toplu ürün senkronizasyonu geliştirdim; değişmeyen kaynaklar atlanıyor,
      kaldırılan ürünler pasife alınıyor.
\item Fiyat değiştiğinde kullanıcıya push bildirimi, firma kaydı ve ödeme sorunlarında ekibe Slack bildirimi gönderen yapıyı ekledim.
\end{itemize}
\vspace{3pt}
{\itshape Software Engineer Intern}\hfill{\color{muted}Haziran 2025 -- Ağustos 2025}\par
\begin{itemize}
\item PostgreSQL indeksleri, Redis önbelleği ve Elasticsearch üzerindeki çalışmalarla ortalama API yanıt süresinin
      \ca{}450~ms'den \ca{}200~ms'ye inmesine katkı sağladım.
\item Mağaza kategorilerini platform kategorileriyle eşleştiren servisi geliştirdim; 160 canlı örneklik test setinde
      \%95 doğru eşleşme elde ettim.
\item Tekrarlanan linkleri temizleyen ve stoğu biten ürünleri arama indeksiyle eşitleyen Django komutları yazdım.
\item Flutter uygulamasının farklı ekran boyutlarına uyarlanmasında çalıştım; 45+ işlev ve arayüz sorununu tespit edip giderdim.
\end{itemize}
\ctx{Proje sonlandırıldı; kaynak kodlar şirketin özel GitHub reposunda bulunuyor.}
\tech{Python\sep Django REST Framework\sep FastAPI\sep PostgreSQL\sep Redis\sep Elasticsearch\sep Celery\sep Flutter}
\end{block}

\begin{block}
\head{Upwork}{Temmuz 2026 -- Halen}
\sub{Software Engineer (Freelance)\sep Uzaktan}
\begin{itemize}
\item Müşteri projelerinde Flutter ile mobil uygulama geliştiriyor, mağaza yayınlarını yürütüyorum.
\end{itemize}
\end{block}

%% ======================= DİĞER PROJELER =======================
\begin{block}[Diğer Projeler]
\head{Anket API\sep\href{https://github.com/enesorhann/Case-Study}{github.com/enesorhann/Case-Study}}{Temmuz 2026}
Anket, soru ve cevap yönetimi sunan .NET 8 API'si (case study). Komutları ve sorguları MediatR ile ayırdım;
anket olaylarını RabbitMQ ile arka plan servisine ilettim.
\end{block}

\begin{block}
\head{Mail Sender Smart Edition\sep\href{https://mailsendersmartedition.com}{mailsendersmartedition.com}}{Ağustos 2025 -- Haziran 2026}
Kullanıcılara kendi alan adımızda e-posta kutusu sunan Django REST API'yi geliştirdim. Gönderim, kurduğum Postfix sunucusundan
SPF, DKIM ve DMARC uyumlu yapılıyor. 65 kayıtlı kullanıcısı olan uygulama
\href{https://apps.apple.com/us/app/mail-sender-smart-edition/id6758946829}{App Store} ve
\mbox{\href{https://play.google.com/store/apps/details?id=com.mailsendersmartedition.enpa}{Google Play}'de} yayında.
\end{block}

\begin{block}
\head{CompAtlas\sep\href{https://compatlas.com.tr/}{compatlas.com.tr}}{Temmuz 2026 -- Halen}
Türkiye'nin 81 ilindeki 17K+ teknoloji firmasının verisini OSM, Google Places ve Overture Maps'ten topladım;
kayıtları temizleyip tekilleştirdim ve LLM ile sınıflandırdım. Platforma kayıt olan firma sayısı artıyor;
Türkiye pazarını araştırmak isteyen yurt dışındaki müşterilerden de destek talepleri geliyor.
\end{block}

%% ======================= YETENEKLER =======================
\begin{block}[Teknik Yetenekler]
\renewcommand{\arraystretch}{1.25}
\noindent\begin{tabular}{@{}>{\raggedright\arraybackslash\bfseries}p{0.22\linewidth}@{}>{\raggedright\arraybackslash}p{0.78\linewidth}@{}}
.NET & C\#, ASP.NET Core Web API, EF Core, SignalR, Hangfire, MediatR, xUnit\\
Python & Django REST Framework, FastAPI, Celery\\
Veri ve mesajlaşma & PostgreSQL, Redis, Elasticsearch, RabbitMQ\\
Mimari ve API & Katmanlı mimari, Repository pattern, CQRS, REST API tasarımı, JWT, RBAC\\
Araçlar ve DevOps & Git, GitHub Actions, Docker, Kubernetes, Linux, Swagger / OpenAPI, Postman\\
Diğer & TypeScript, Next.js, Flutter, Firebase\\
\end{tabular}
\end{block}

%% ======================= EĞİTİM =======================
\begin{block}[Eğitim]
\head{Konya Teknik Üniversitesi}{Ağustos 2022 -- Haziran 2026}
Bilgisayar Mühendisliği (Lisans)\sep Genel not ortalaması 3.39 / 4.00\par
\vspace{3pt}
Sertifikalar:
\href{https://drive.google.com/file/d/1btYql9fKCvBwqHeh06F_NE257FhceNCB/view}{Sertifika portföyü}\sep
\href{https://drive.google.com/file/d/12AcWLnSGEfcZUWUb945n6nD-iPB7aIA9/view?usp=sharing}{Teknofest Katılım Belgesi}\par
\vspace{3pt}
Referans: Doç. Dr. Sait Ali Uymaz, Konya Teknik Üniversitesi Bilgisayar Mühendisliği
\end{block}

\end{document}
